\chapter{ภาคผนวกแรก (The first appendix)}

ภาคผนวกนี้รวบรวมเอกสารอ้างอิงเสริม ข้อมูลเชิงเทคนิค (Technical Details) และเครื่องมือปฏิบัติการที่เกี่ยวข้องกับกระบวนการพัฒนาระบบแพลตฟอร์มพาณิชย์อิเล็กทรอนิกส์ "Elder Store" ตลอดจนระบบแผงควบคุมหลังบ้าน (Admin Dashboard) สำหรับผู้ประกอบการร้านค้า โดยประกอบไปด้วยเครื่องมือแบบสอบถามที่ใช้ในการเก็บรวบรวมข้อมูลพฤติกรรมผู้ใช้งาน (User Research Instrument) โครงสร้างชุดคำสั่งอัลกอริทึมบางส่วนที่เป็นแกนหลักของระบบ และคำอธิบายเชิงลึกของสถาปัตยกรรมฐานข้อมูลเชิงสัมพันธ์ \cite{postgresql2023} ที่คอยขับเคลื่อนโครงงานนี้

\section{แบบสำรวจความต้องการและพฤติกรรมการใช้งาน (User Research Survey)}
แบบสอบถามชุดนี้ถูกจัดทำขึ้นตามหลักการวิศวกรรมความต้องการ (Requirements Engineering) \cite{pressman2014} เพื่อใช้เป็นเครื่องมือในการเก็บรวบรวมข้อมูลปฐมภูมิ (Primary Data) เจาะจงไปที่กลุ่มประชากรผู้สูงอายุและผู้ดูแล เพื่อนำผลลัพธ์ที่ได้มาสังเคราะห์เป็นข้อกำหนดของระบบ (System Specifications) เผยให้เห็นถึงข้อจำกัดและคอขวดในการใช้งานเทคโนโลยี (Pain Points) \cite{wai_older_users} และนำไปสู่การออกแบบแพลตฟอร์มที่สะท้อนถึงหลักอารยสถาปัตย์ (Inclusive Design) \cite{inclusive2022} อย่างแท้จริง

\textbf{ส่วนที่ 1: พฤติกรรมการบริโภคและการเข้าถึงเทคโนโลยี}
\begin{enumerate}
  \item ปัจจุบันท่านมีการสั่งซื้อสินค้าอุปโภคบริโภค หรือสินค้าเพื่อสุขภาพผ่านช่องทางออนไลน์บ่อยเพียงใด? (เป็นประจำ / นานๆ ครั้ง / ไม่เคยเลย)
  \item อุปกรณ์ที่ท่านรู้สึกถนัดและใช้งานบ่อยที่สุดในการเข้าถึงอินเทอร์เน็ตคืออุปกรณ์ใด? (โทรศัพท์สมาร์ตโฟน / คอมพิวเตอร์ตั้งโต๊ะ / แท็บเล็ต) 
\end{enumerate}

\textbf{ส่วนที่ 2: อุปสรรคและข้อจำกัดในการใช้งานแพลตฟอร์มออนไลน์ในปัจจุบัน}
\begin{enumerate}
    \setcounter{enumi}{2}
    \item ปัญหาหลักที่ท่านมักพบเมื่อพยายามซื้อสินค้าผ่านเว็บไซต์หรือแอปพลิเคชันรูปแบบเดิมคืออะไร? (เลือกตอบได้มากกว่า 1 ข้อ) \cite{accessibility_elderly_2012}
    \begin{itemize}
        \item ตัวอักษรและปุ่มกดมีขนาดเล็กเกินไป ส่งผลให้มองเห็นไม่ชัดเจน หรือนิ้วสัมผัสคลาดเคลื่อนเป้าหมาย
        \item สีสันของแพลตฟอร์มมีความฉูดฉาดและกลืนกัน ทำให้แยกแยะองค์ประกอบเมนูได้ลำบาก
        \item ขั้นตอนการเลือกซื้อหรือการพยายามเก็บคูปองส่วนลดมีความซับซ้อนและมีเมนูย่อยมากเกินไป
        \item ความกังวลเรื่องความปลอดภัย กลัวข้อมูลส่วนตัวรั่วไหล หรือกลัวถูกกระบวนการมิจฉาชีพหลอกลวง
    \end{itemize}
    \item ในขั้นตอนการทำธุรกรรม ท่านคิดว่าการชำระเงินผ่านระบบสแกนคิวอาร์โค้ดตามมาตรฐานของรัฐ (PromptPay) สะดวกและตอบโจทย์ความปลอดภัยของท่านมากกว่าการพิมพ์รหัสบัตรเครดิตลงในหน้าเว็บโดยตรงหรือไม่? \cite{promptpay2021}
\end{enumerate}

\textbf{ส่วนที่ 3: ความต้องการฟังก์ชันเสริมเพื่อยกระดับประสบการณ์ใช้งาน}
\begin{enumerate}
    \setcounter{enumi}{4}
    \item ฟีเจอร์หรือระบบจัดการใดต่อไปนี้ ที่ท่านคิดว่าจะช่วยให้การซื้อของออนไลน์สำหรับท่านง่ายขึ้น? (เลือกตอบได้มากกว่า 1 ข้อ)
    \begin{itemize}
        \item ขนาดตัวอักษรตั้งต้นบนเบราว์เซอร์ที่นำเสนอแบบขยายใหญ่เป็นพิเศษ (Large Typography)
        \item ระบบการใช้อัลกอริทึมแนะนำผลิตภัณฑ์ที่สอดคล้องกับหมวดหมู่ความสนใจของผู้สูงอายุอัตโนมัติ
        \item กระบวนการชำระเงินแบบรวบรัดขั้นตอนให้เหลือเพียงหน้าจอเดียว (Single-step Checkout) \cite{uxtrends2023}
        \item ปุ่มสำหรับคอลเซ็นเตอร์ติดต่อผู้ดูแลระบบ (Admin Support) ที่สามารถป๊อปอัปขึ้นมาได้ทันทีเมื่อเกิดข้อสงสัย
    \end{itemize}
\end{enumerate}

\section{สถาปัตยกรรมความปลอดภัยและนโยบายคุ้มครองข้อมูลส่วนบุคคล (Security Architecture \& PDPA Compliance)}

เนื่องจากแพลตฟอร์มพาณิชย์อิเล็กทรอนิกส์ "Elder Store" เป็นระบบที่มีศูนย์กลางขับเคลื่อนด้วยการทำธุรกรรมทางการเงินและมีการจัดเก็บที่อยู่ผู้ใช้งาน ซึ่งบรรจุองค์ประกอบของข้อมูลส่วนบุคคลที่มีความอ่อนไหว (Sensitive Data) โดยเฉพาะประชากรกลุ่มผู้สูงอายุที่อาจมีความเสี่ยงต่อการตกเป็นเหยื่ออาชญากรรมทางวิศวกรรมสังคม (Social Engineering) ได้ง่าย คณะผู้จัดทำจึงได้ยกระดับความสำคัญของการออกแบบสถาปัตยกรรมความมั่นคงปลอดภัยของระบบ (Security Architecture) ภายใต้กรอบมาตรฐานวิศวกรรมซอฟต์แวร์ และพระราชบัญญัติคุ้มครองข้อมูลส่วนบุคคล พ.ศ. 2562 (PDPA) ดังมีรายละเอียดเชิงลึกต่อไปนี้:

\subsection{1. มาตรการเข้ารหัสและพิสูจน์ยืนยันตัวตน (Authentication and Cryptography)}
ระบบได้นำโปรโตคอลการพิสูจน์ตัวตนแบบไร้สถานะ (Stateless Authentication) ผ่านมาตรฐานอุตสาหกรรม JSON Web Token (JWT) มาประยุกต์ใช้เป็นคอร์หลักในการจัดการช่วงเวลาเซสชัน (Session Management) โดยเมื่อผู้สูงอายุหรือแอดมินทำการเข้าสู่ระบบด้วยพาสเวิร์ดสำเร็จ เซิร์ฟเวอร์จะทำการผลิตและส่งมอบโทเค็น (Token) ที่ถูกเข้ารหัสด้วยลายเซ็นดิจิทัล (Digital Signature) เพื่อใช้เป็นใบเบิกทางในการเข้าถึงเส้นทาง API (API Routes) ต่างๆ 

ในทวีปของการจัดเก็บรหัสผ่าน (Password Storage) ลงสู่ระบบเลเยอร์ฐานข้อมูล PostgreSQL คณะผู้จัดทำได้วางสคริปต์บังคับใช้ฟังก์ชันการแฮชข้อมูลแบบทิศทางเดียว (One-way Hashing Function) ผ่านอัลกอริทึมที่มีความต้านทานระดับสูง (อาทิ Bcrypt) ควบคู่กับการสุ่มเติมค่าเกลือ (Salting Pattern) ซึ่งจะแปรสภาพรหัสผ่านดั้งเดิมของผู้ใช้ให้กลายเป็นสายอักขระตัวถอดรหัสที่มนุษย์ไม่สามารถอ่านไขได้ (Unreadable Ciphertext) ส่งผลให้แม้ประการผู้ดูแลระบบหลังบ้าน (Database Administrator) หรือทีมนักพัฒนา ก็ไม่สามารถล่วงรู้และมองเห็นรหัสผ่านที่แท้จริงของลูกค้าประดุจแนวคิด Zero-Knowledge Privacy ยกเว้นแต่ตัวลูกค้าเอง

\subsection{2. การปกป้องข้อมูลระหว่างรอยต่อการส่งผ่าน (Data in Transit Protection)}
เพื่อป้องกันกลุ่มมิจฉาชีพที่มีพฤติกรรมดักจับแพ็กเกจข้อมูลกลางอากาศ (Man-in-the-Middle Attack) ทราฟฟิกข้อมูลที่สัญจรไปมาระหว่างหน้าเบราว์เซอร์ของผู้ใช้งาน (Client-side) และฝั่งเครื่องรันไทม์แม่ข่าย (Server-side) ทั้งหมด จะถูกบีบอัดและบังคับให้ส่งพาดผ่านช่องทางอุโมงค์สถาปัตยกรรมโปรโตคอล HTTPS (Hypertext Transfer Protocol Secure) ซึ่งถูกครอบเกราะคุ้มครองด้วยมาตรฐานการเข้ารหัส SSL/TLS (Secure Sockets Layer/Transport Layer Security) ระดับ 256-bit แบบปลายทางถึงปลายทาง (End-to-End Encryption) ระบบจึงการันตีได้ว่าข้อมูลบิลการสั่งซื้อ ที่อยู่พิกัดในการจัดส่งพัสดุ และไฟล์รูปภาพสลิปที่โอนชำระเงิน จะได้รับการคุ้มครองจากการถูกดักแทรกแซงโดยสถาบันองค์กรแฮกเกอร์ภายนอกเครือข่าย

\subsection{3. การดำเนินงานประยุกต์ตามเจตนารมณ์กฎหมาย PDPA (Privacy By Design)}
คณะผู้จัดทำบูรณาการกระบวนทัศน์แนวคิด "ความเป็นส่วนตัวตั้งแต่ขั้นตอนการออกแบบ" (Privacy by Design) ผูกติดไว้เป็นแกนกลางหลักของแพลตฟอร์ม Elder Store โดยมีมาตรการเชิงประจักษ์ที่นำไปปฏิบัติจริง (Implementation) ดังนี้:

\begin{itemize}
    \item \textbf{ระบบขออนุญาตความยินยอมป๊อปอัปคุกกี้ (Cookie Consent Banner):} ทันทีที่มีเส้นทราฟฟิกผู้ใช้งานเปิดเข้ามาเหยียบบนอินเทอร์เฟซของระบบ แพลตฟอร์มจะแสดงชุดหน้าต่างแบนเนอร์เชิงซ้อนเพื่อร้องขอสิทธิ์ (Consent) ในการจัดเก็บไฟล์ข้อมูลคุกกี้ขนาดย่อมลงบนเครื่องฮาร์ดแวร์สมาร์ตโฟนของผู้ใช้ โดยมอบอำนาจเด็ดขาดให้บัญชีเลือกว่าจะ "ยอมรับ" หรือ "ปฏิเสธ" การถูกติดตามพฤติกรรมนี้ ซึ่งเป็นกุญแจสำคัญในการรักษาจริยธรรมความโปร่งใส
    \item \textbf{หลักวิศวกรรมการทอนข้อมูลเท่าที่จำเป็น (Data Minimization):} ในฟอร์มสมัครสมาชิกระยะเริ่มต้น และหน้าต่างดำเนินการชำระเงิน ระบบจะเรียกร้องค่าบังคับคีย์เข้าเฉพาะจุดพิกัดข้อมูลที่ต้องนำไปคำนวณการจัดส่งพัสดุและระบุตัวตนทางธุรกรรมอย่างหลีกเลี่ยงไม่ได้ (ประกอบด้วย ตำแหน่งชื่อจริง, เบอร์โทรศัพท์, และรหัสไปรษณีย์) โดยจะปิดกั้นฟิลด์คำถามที่ไม่จำเป็น เช่น เลขบัตรประจำตัวประชาชน 13 หลัก วันเกิด หรือฐานรายได้ เพื่อตีวงจำกัดกรอบความสูญเสียไม่ให้ลุกลามหากเกิดเหตุการณ์เซิร์ฟเวอร์โดนเจาะทะลวงผ่านสถานการณ์ที่เลวร้ายที่สุด (Worst-case Scenario)
    \item \textbf{สิทธิการเข้าควบคุมสัมปทานข้อมูลส่วนบุคคล (Right to Modify \& Erasure):} แพลตฟอร์มกระจายอำนาจส่งกลับไปยังเจ้าของบัญชีผู้สูงอายุ โดยเมื่อพวกเขาทำการยืนยันตัวตนสำเร็จ จะมีสิทธิแท็บเข้าแก้ไขอัปเดตที่อยู่ตนเองในฐานระบบ (Data Rectification) หรือแม้กระทั่งพาดขอสิทธิลบทำลายประวัติบัญชีตนเองผ่านเซิร์ฟเวอร์แบบเบ็ดเสร็จ (Hard-delete mechanism) อันมีความสอดรับกับปรัชญากฎหมาย "สิทธิที่จะถูกลืม" (Right to be Forgotten) ตามองค์ประกอบบริบทพระราชบัญญัติยุคความมั่นคงไซเบอร์ใหม่ทุกประการ
\end{itemize}



\section{พจนานุกรมข้อมูลและรายละเอียดโครงสร้างฐานข้อมูล (Database Dictionary)}
แพลตฟอร์ม Elder Store ใช้ระบบจัดการฐานข้อมูลเชิงสัมพันธ์ระดับองค์กร (PostgreSQL) \cite{postgresql2023} ในการรวบรวมและประมวลผลขุมข้อมูล ซึ่งได้รับการออกแบบโครงสร้างตารางข้อมูล (Table Schema) ให้มีความเป็นบรรทัดฐานขั้นสูงระดับที่ 3 (Third Normal Form: 3NF) เพื่อป้องกันการเกิดความผิดปกติรบกวนข้อมูลเวลาย้ายข้อมูล (Data Anomalies) การลดความซ้ำซ้อน และรองรับอัตราการเติบโตทางทรัพยากรผู้ใช้ในอนาคต ดังรายละเอียดแบบจำลองระดับตรรกะต่อไปนี้:

% =========================================================
% ตารางที่ 1: Users
% =========================================================
\subsection{ตารางข้อมูลผู้ใช้งาน (tb\_users)}
ใช้สำหรับจัดเก็บข้อมูลคลังบัญชีผู้ใช้งานส่วนตัวทั้งของผู้ซื้อ (ผู้สูงอายุ) และผู้ดูแลระบบ (Admin)
\begin{table}[h!]
\centering
\caption{พจนานุกรมข้อมูลตาราง tb\_users}
\vspace{0.2cm}
\begin{tabular}{|l|l|c|l|}
\hline
\textbf{ชื่อฟิลด์ (Field Name)} & \textbf{ชนิดข้อมูล (Type)} & \textbf{คีย์ (Key)} & \textbf{คำอธิบาย (Description)} \\ \hline
user\_id & UUID & PK & รหัสอ้างอิงประจำตัวผู้ใช้งานบัญชี (สร้างอัตโนมัติ) \\ \hline
email & VARCHAR(255) & UK & อีเมลเข้าสู่ระบบ (ระบบไม่อนุญาตให้ซ้ำกัน) \\ \hline
password\_hash & VARCHAR(255) & - & รหัสผ่านที่ผ่านการรักษาความปลอดภัยแบบแฮช \\ \hline
first\_name & VARCHAR(100) & - & ชื่อจริงระบุตัวตนของผู้ใช้งานระบบ \\ \hline
last\_name & VARCHAR(100) & - & นามสกุลจริงของผู้ใช้งานระบบ \\ \hline
phone\_number & VARCHAR(15) & - & เบอร์โทรศัพท์สำหรับใช้ระบุตัวจัดส่งพัสดุ \\ \hline
role & VARCHAR(20) & - & กำหนดมาตรการระดับสิทธิ์ (เช่น 'user', 'admin') \\ \hline
created\_at & TIMESTAMP & - & ตราประทับเวลา (Timestamp) ของการสมัครบัญชีสำเร็จ \\ \hline
\end{tabular}
\end{table}

% =========================================================
% ตารางที่ 2: Products
% =========================================================
\subsection{ตารางข้อมูลสินค้า (tb\_products)}
ศูนย์รวมข้อมูลรายการสินค้าทั้งหมดที่ถูกนำมาประกาศจัดจำหน่ายบนหน้าแดชบอร์ด
\begin{table}[h!]
\centering
\caption{พจนานุกรมข้อมูลตาราง tb\_products}
\vspace{0.2cm}
\begin{tabular}{|l|l|c|l|}
\hline
\textbf{ชื่อฟิลด์ (Field Name)} & \textbf{ชนิดข้อมูล (Type)} & \textbf{คีย์ (Key)} & \textbf{คำอธิบาย (Description)} \\ \hline
product\_id & UUID & PK & รหัสอ้างอิงโครงสร้างประจำตัวสินค้า \\ \hline
product\_name & VARCHAR(255) & - & ชื่อเรียกหรือยี่ห้อพาดหัวของสินค้า \\ \hline
description & TEXT & - & คำบรรยายสรรพคุณหรือรายละเอียดสินค้าครบถ้วน \\ \hline
price & DECIMAL(10,2) & - & ต้นทุนค่าตัวสินค้าต่อหน่วย (รูปแบบทศนิยมสกุลบาท) \\ \hline
stock\_quantity & INT & - & คำนวณจำนวนไอเทมที่เหลืออยู่พร้อมจัดส่งในดรอปชิป \\ \hline
image\_url & VARCHAR(512) & - & ลิงก์เก็บพาร์ทที่ตั้งไฟล์รูปภาพตัวอย่างแบนเนอร์สินค้า \\ \hline
image\_alt\_text & VARCHAR(255) & - & คำอธิบายโครงภาพสำหรับซอฟต์แวร์อ่านหน้าจอ \cite{wcag2023} \\ \hline
is\_active & BOOLEAN & - & สถานะการควบคุมตรรกะบูลีนการแสดงผลทางหน้าเว็บ \\ \hline
\end{tabular}
\end{table}

% =========================================================
% ตารางที่ 3: Orders
% =========================================================
\subsection{ตารางข้อมูลคำสั่งซื้อ (tb\_orders)}
ตารางหลักตัวกลางที่บันทึกร่องรอยเส้นทางการทำธุรกรรมพาณิชย์และข้อมูลเครือข่ายจุดหมายปลายทาง
\begin{table}[h!]
\centering
\caption{พจนานุกรมข้อมูลตาราง tb\_orders}
\vspace{0.2cm}
\begin{tabular}{|l|l|c|l|}
\hline
\textbf{ชื่อฟิลด์ (Field Name)} & \textbf{ชนิดข้อมูล (Type)} & \textbf{คีย์ (Key)} & \textbf{คำอธิบาย (Description)} \\ \hline
order\_id & UUID & PK & รหัสหมายเลขใบสั่งซื้อ (บิลเรียกเก็บเงิน) \\ \hline
user\_id & UUID & FK & รหัสผู้ใช้ที่กระทำการเชื่อมจากตาราง tb\_users \\ \hline
total\_amount & DECIMAL(10,2) & - & ยอดชำระสุทธิบวกรวมภาษีของคำสั่งซื้อบิลนี้ \\ \hline
order\_status & VARCHAR(50) & - & แสดงสถานะวงจรพัสดุ (Pending, Shipped, Cancelled) \\ \hline
shipping\_address & TEXT & - & ข้อมูลปักหมุดสถานที่จัดส่งพัสดุดิจิทัลฉบับสมบูรณ์ \\ \hline
tracking\_number & VARCHAR(100) & - & หมายเลขคาร์โก้ติดตามพัสดุร่วมกับบริษัทขนส่ง \\ \hline
created\_at & TIMESTAMP & - & วันและเวลาสแตมป์ที่ลูกค้ากดยืนยันคำสั่งชำระเงิน \\ \hline
\end{tabular}
\end{table}

% =========================================================
% ตารางที่ 4: Order Items
% =========================================================
\subsection{ตารางข้อมูลรายการสินค้าในคำสั่งซื้อ (tb\_order\_items)}
ตารางโครงสร้างลูกของคำสั่งซื้อ สำหรับทำฟังก์ชันแตกรายการย่อย กรณีซื้อของหลายชนิดต่อครั้ง
\begin{table}[h!]
\centering
\caption{พจนานุกรมข้อมูลตาราง tb\_order\_items}
\vspace{0.2cm}
\begin{tabular}{|l|l|c|l|}
\hline
\textbf{ชื่อฟิลด์ (Field Name)} & \textbf{ชนิดข้อมูล (Type)} & \textbf{คีย์ (Key)} & \textbf{คำอธิบาย (Description)} \\ \hline
item\_id & UUID & PK & รหัสตัวกำหนดบันทึกส่วนย่อยของตะกร้า \\ \hline
order\_id & UUID & FK & อาศัยการเชื่อมสะพานคีย์หลักมาจาก tb\_orders \\ \hline
product\_id & UUID & FK & เชื่อมตารางเพื่อดึงออบเจ็กต์ข้อมูลจากส่วน tb\_products \\ \hline
quantity & INT & - & จำนวนชิ้นคูณทวีคูณของสินค้าประเภทนี้ที่จะต้องส่งออก \\ \hline
unit\_price & DECIMAL(10,2) & - & รักษาราคาสินค้า ณ วันเวลาที่ทำการบันทึก (กันราคาเหวี่ยง) \\ \hline
\end{tabular}
\end{table}

% =========================================================
% ตารางที่ 5: Payments
% =========================================================
\subsection{ตารางข้อมูลการชำระเงิน (tb\_payments)}
คอยควบคุมรักษาสถานะและดูแลสิทธิหลักฐานความโปร่งใสทางธุรกรรม (Proof of Payment) ให้ผู้บริโภค
\begin{table}[h!]
\centering
\caption{พจนานุกรมข้อมูลตาราง tb\_payments}
\vspace{0.2cm}
\begin{tabular}{|l|l|c|l|}
\hline
\textbf{ชื่อฟิลด์ (Field Name)} & \textbf{ชนิดข้อมูล (Type)} & \textbf{คีย์ (Key)} & \textbf{คำอธิบาย (Description)} \\ \hline
payment\_id & UUID & PK & รหัสด้านบัญชีที่แยกต่างหากสำหรับการอนุมัติเงินแต่ละงวด \\ \hline
order\_id & UUID & FK & ปฏิสัมพันธ์เพื่ออ้างอิงยืนยันการจ่ายว่าผูกกับ tb\_orders ใด \\ \hline
payment\_method & VARCHAR(50) & - & ระบบช่องทางชำระเงินขาเข้า (เช่น PromptPay, Credit Card) \\ \hline
payment\_slip\_url & VARCHAR(512) & - & แนบสตริงลิงก์รูปภาพสลิปที่วิ่งเข้าไปเก็บสู่ระบบเซอร์เวอร์บนคลาวด์ \\ \hline
payment\_status & VARCHAR(50) & - & สถานะสลิปจากเงื่อนไขระบบ (เช่น รอการตรวจสอบ, อนุมัติแล้ว) \\ \hline
verified\_at & TIMESTAMP & - & วินาทีที่ผู้ดูแลระบบหลังบ้านตอบรับการอนุมัติยอดเงินเข้าสู่ตัวเลขจริง \\ \hline
\end{tabular}
\end{table}

\chapter{ภาคผนวก ข (The second appendix)}

ภาคผนวกนี้จัดทำขึ้นเพื่อเปรียบเสมือนคู่มือแนะแนวทางการใช้งาน (User Manual) แพลตฟอร์มพาณิชย์อิเล็กทรอนิกส์ "Elder Store" อย่างเป็นลำดับขั้นตอน โดยแบ่งออกเป็น 2 ส่วนหลัก ได้แก่ คู่มือการใช้งานส่วนหน้าสำหรับผู้สูงอายุ (Frontend) ที่เน้นกระบวนการสั่งซื้อให้สอดคล้องกับพฤติกรรมศาสตร์ และคู่มือสำหรับผู้ดูแลระบบ (Admin Dashboard) เพื่อใช้ในการบริหารจัดการฐานข้อมูลคำสั่งซื้อและดูแลความรัดกุมของระบบ

\section{คู่มือการใช้งานสำหรับผู้สูงอายุ (Frontend User Guide)}
ระบบส่วนหน้าถูกออกแบบมาโดยการตัดทอนเมนูที่ซับซ้อนทิ้งไป ยึดถือหลักการ "อย่าทำให้ผู้ใช้งานต้องคิดกระบวนการขั้นตอนด้วยตนเอง" ส่งผลให้คู่มือการใช้งานมีลำดับการดำเนินการที่กระชับ ดังนี้:

\subsection{1. ขั้นตอนการลงทะเบียนและการเข้าสู่ระบบ}
\begin{enumerate}
    \item เมื่อเข้าสู่หน้าแรกของเว็บไซต์ผ่านเบราว์เซอร์ ให้ผู้ใช้นิ้วคลิกหรือสัมผัสที่ปุ่ม \textbf{"เข้าสู่ระบบ / สมัครสมาชิก"} บริเวณมุมขวาบนของหน้าจอ
    \item หากเข้าใช้งานเป็นครั้งแรก ให้เลือกแท็บ \textbf{"สร้างบัญชีใหม่"} จากนั้นพิมพ์กรอกข้อมูลอัตลักษณ์ส่วนบุคคลเบื้องต้น เช่น ชื่อ เบอร์โทรศัพท์ และรหัสผ่านแบบสั้น
    \item เมื่อระบบประมวลผลสมัครสำเร็จ มันจะนำพาท่านส่งกลับเข้าสู่หน้าโฮมเพจหลักในสถานะที่ระบบจดจำตัวตน (ล็อกอิน) เรียบร้อยแล้ว
\end{enumerate}

\subsection{2. ขั้นตอนการค้นหาและหยิบสินค้า}
\begin{enumerate}
    \item ผู้ใช้งานสามารถกวาดสายตาเลื่อนดูสินค้าแนะนำ (Recommended Items) บนหน้าแรก หรือเลือกหมวดหมู่ย่อยที่ต้องการจากแถบเมนูนำทาง (Navigation Bar)
    \item ทัชสกรีนที่รูปภาพแบนเนอร์สินค้าชิ้นนั้นๆ เพื่อเข้าสู่ \textbf{"หน้ารายละเอียดสินค้า"} ซึ่งจะบรรยายสรรพคุณ รีวิว และประทับตราป้ายราคาด้วยตัวอักษรขนาดใหญ่
    \item หากคุณตัดสินใจสั่งซื้อ ให้กดที่ปุ่ม \textbf{"เพิ่มลงตะกร้า"} สีส้มขนาดใหญ่ (Touch Target size $> 48$px) บริเวณฐานล่างสุดของหน้าจอ
\end{enumerate}

\subsection{3. ขั้นตอนการสรุปตะกร้าสินค้าและการชำระเงิน (Checkout Flow)}
\begin{enumerate}
    \item แตะที่ไอคอน \textbf{ตะกร้าสินค้า} บริเวณมุมขวาบน เพื่อรีเช็กรายการสินค้าทั้งหมดที่เตรียมจะชำระเงิน ระบบจะคำนวณราคาสุทธิให้อัตโนมัติ
    \item กรอกที่อยู่บ้านเพื่อการจัดส่งพัสดุ หากเคยซื้อแล้ว ระบบจะดึงข้อมูลที่อยู่เดิมมาเตรียมรอไว้ หากตรวจสอบถูกต้องให้กดปุ่ม \textbf{"ดำเนินการชำระเงิน"}
    \item ระบบจะสร้างรหัส \textbf{QR Code (PromptPay)} ส่วนกลางของแพลตฟอร์ม ให้ลูกค้านำแอปพลิเคชันธนาคารในมือถือสแกนชำระเงินยอดตรงตามจริง
    \item ลำดับสุดท้าย กดปุ่ม \textbf{"แนบหลักฐานการโอนเงิน"} และเลือกรูปสลิปจากคลังสมุดภาพในโทรศัพท์ ก่อนกด \textbf{"ยืนยันคำสั่งซื้อ"} อย่างสมบูรณ์
\end{enumerate}

\subsection{4. ขั้นตอนการทวงถามหรือติดตามสถานะพัสดุ}
\begin{enumerate}
    \item ให้ผู้ใช้กดเข้าไปที่ไอคอนโปรไฟล์ แล้วเข้าเมนู \textbf{"ประวัติการซื้อของฉัน"}
    \item โปรแกรมจะลิสต์สถานะตั๋วแบบเรียลไทม์ เช่น "กำลังตรวจสอบสลิป", "กำลังจัดส่งออก", "สำเร็จแล้ว" 
    \item ในสถานะกำลังจัดส่ง ผู้สูงอายุจะเห็นตัวเลขรหัสติดตามพัสดุ (Tracking No.) ที่แอดมินพิมพ์ฝากไว้ให้ตรวจสอบร่วมกับบริษัทไปรษณีย์
\end{enumerate}

\section{คู่มือการใช้งานฐานบัญชาการสำหรับผู้ดูแลระบบ (Admin Guide)}
การบริหารระบบหลังบ้าน (Dashboard) ถูกสร้างขึ้นเพื่อติดตามวิเคราะห์สถิติ ทรัพยากรและแก้ไขช่องโหว่ โดยจะสงวนกลไกให้เปิดได้เฉพาะผู้ใช้ที่มีเลเวลสิทธิ์ \textbf{'admin'} ในระบบฐานข้อมูลเท่านั้น:

\subsection{1. ขั้นตอนการพิสูจน์เพื่อเข้าสู่ระบบหลังบ้าน}
\begin{enumerate}
    \item พิมพ์เข้าถึงผ่านเส้นทาง URL ไพรเวท (เช่น \texttt{/admin-login}) ระบบเครือข่ายจะดักจับและบังคับให้ยืนยันตัวตนด้วยบัญชีอีเมลระดับผู้ดูแลที่ได้รับอนุญาต
    \item เมื่อเข้าสู่ระบบสำเร็จ ตัวเซิร์ฟเวอร์จะคืนค่าหน้าต่างควบคุมหลัก (Dashboard Overview Analytics) ให้คุณตรวจสอบอัตราการเจริญเติบโตของยอดขายรายวัน 
\end{enumerate}

\subsection{2. ขั้นตอนการจัดการแคตตาล็อกสินค้าออนไลน์ (Product Management)}
\begin{enumerate}
    \item เลื่อนปุ่มไปที่เมนู \textbf{"จัดการสินค้า"} ทางแถบคอนโซลฝั่งซ้ายมือ
    \item กดไอคอน \textbf{"+ เพิ่มสินค้าใหม่"} คุณจะต้องกรอกข้อมูลตัวแปรชื่อ, ราคา, อัปโหลดโฮสต์รูปภาพ, จำนวนสต็อก และที่สำคัญคือคำอธิบายสื่อภาพ (Alt-text)
    \item หากมีสินค้าตัวใดที่หมดสต็อกระยะยาว หรือผู้ประกอบการต้องการพักการขาย คุณมีสิทธิกดสวิตช์เปิด-ปิดทางบูลีน (Status Toggle) เพื่อซ่อนการแสดงผลหน้าบ้านได้ทันที
\end{enumerate}

\subsection{3. ขั้นตอนการตรวจสอบและตอบรับสถานะคำสั่งซื้อ (Order Reviewing)}
\begin{enumerate}
    \item แอดมินมีหน้าที่ความรับผิดชอบต้องคอยสแตนด์บายประเมินรายงานสั่งซื้อใหม่ที่ไหลเข้ามาในเมนู \textbf{"คำสั่งซื้อทั้งหมด"}
    \item คลิกเลือกบิลที่ค้างสถานะ \textit{"รอตรวจสอบยอดโอน"} เพื่อดาวน์โหลดและตรวจสอบรูปภาพหลักฐานสลิปเงินที่ลูกค้าผู้สูงอายุโอนแนบเข้าสู่ระบบคลาวด์
    \item หากยอดเงินในบัญชีบริษัทเข้าและตรงปก ให้กดยืนยันอัปเดตสเตทออเดอร์ให้เป็น \textbf{"ชำระเงินแล้ว"} และหลังจากกระบวนการแพ็คของส่งไปรษณีย์เสร็จ คุณควรนำหมายเลขพัสดุมากรอกอัปเดตลงในระบบเพื่อให้ลูกค้าดูอัปเดต 
\end{enumerate}

\subsection{4. ขั้นตอนการจัดการคูปองส่งเสริมการตลาด (Coupon System)}
\begin{enumerate}
    \item สลับมาที่เมนู \textbf{"จัดการคูปองส่วนลด"} แอดมินสามารถกดบวกเพื่อเจนเนอเรทโค้ดสตริงส่วนลดใหม่เข้าสู่กระแสตลาด
    \item คุณมีสิทธิเซตตัวแปรเงื่อนไขต่างๆ เช่น ปริมาณเปอร์เซ็นต์ส่วนลด, จำนวนเงื่อนไขเงินขั้นต่ำ, จำนวนโค้ดที่แจก และจำกัดวันหมดอายุอัลกอริทึม
    \item โค้ดที่แอกทิฟ (Active) แล้ว จะพร้อมกระจายให้ผู้ใช้งานนำไปกรอกใช้ฟาดฟันลดยอดบิลในหน้าตะกร้าสินค้าหน้าบ้านได้ 
\end{enumerate}

\vspace{1.5cm}
\begin{center}
    \textit{หมายเหตุ: สามารถตรวจสอบภาพประกอบหน้าจออินเทอร์เฟซจริงแบบเต็มความละเอียด ของคู่มือปฏิบัติการแต่ละขั้นตอน ได้ใน เนื้อหาหลัก บทที่ 4 การทดลองและผลลัพธ์การพัฒนาระบบ}
\end{center}
